\chapter{AI 使用情况与人机协作}
\section{两类 AI 使用场景}
本项目需要区分运行时 AI 与开发辅助 AI。运行时 AI 是顾客实际使用的文字、图片或语音入口，影响接口耗时、费用和业务正确性；开发辅助 AI 用于解释问题、辅助编码或文档整理，影响开发效率和成员对成果的理解。两者不能以同一个“使用了 AI”概括，也不能仅凭系统配置推断开发人员使用过某款编码工具。
\section{运行时 AI 的名称与用途}
下表依据仓库 \code{elm_bk/src/main/resources/application.yml} 以及部署配置整理。默认模型名代表代码配置，不表示本次报告整理时已经对外部服务做了真实调用或性能测试。
\begin{longtable}{L{.25\linewidth}L{.27\linewidth}L{.39\linewidth}}
\caption{系统 AI 能力及配置依据}\\\toprule AI 服务/模型&使用场合&作用与证据边界\\\midrule\endfirsthead\toprule AI 服务/模型&使用场合&作用与证据边界\\\midrule\endhead
DeepSeek / \code{deepseek-chat}&文字问答和辅助点餐&在本地意图分流、真实业务数据和权限约束下按需组织回答。配置有总超时与重试；外部可用性需真实联调。\\
DashScope / \code{qwen-vl-plus}&图片识别入口&部署文档允许配置视觉服务；默认模型名可由环境变量覆盖。图片结果仍需转化并核对业务条件。\\
DashScope / \code{qwen3-asr-flash}&语音识别入口&把语音输入转为可继续处理的信息；浏览器权限、模型配置与服务额度会影响可用性。\\
本地规则与业务查询&规则咨询、本人订单查询和失败回退&直接读取服务端业务事实，不必为每次请求调用语言模型；本地规则本身不属于外部大模型。\\\bottomrule
\end{longtable}
外部模型密钥由本地或部署环境提供，不写入报告或仓库。运行时 AI 的具体代码实现见第四章；调用次数、重试放大和并发问题的分析及改进方案见第九章。
\section{开发与文档辅助 AI}
本次报告的材料整理、代码与 Git 记录核对、LaTeX 编写、编译和版面检查使用 Codex 辅助完成。成员提供真实分工、答辩问题和模板保留要求，决定内容范围并纠正初稿偏差。

项目代码开发阶段每位成员实际使用的 AI 工具名称、使用次数及具体采用片段，现有仓库不足以独立证明，仍需成员补充。不能将运行时使用 DeepSeek 推导为开发时使用 DeepSeek，也不能从提交风格断言使用过 ChatGPT、Cursor 或其他工具。

\section{人机协作的工作过程}
本次报告整理采用以下可核对流程：首先由成员提供 SRS、项目仓库、分工、老师意见与交叉测试材料，确定事实边界；随后借助 Codex 检索代码和提交记录，整理章节、生成 LaTeX 并编译；成员指出原 SRS 和模板必须完整保留后，调整为原始文件逐字节保留、通过独立入口追加新章节；最后核对原件一致性、材料来源、统计口径和页面呈现。这个过程说明，人提出需求和纠正偏差，AI 承担整理与执行，最终内容仍需人理解和确认。

对开发中的人机协作，应由成员先给出需求、约束、输入输出和测试标准，AI 再给出候选实现；成员阅读差异、运行测试并检查边界后决定是否采纳。没有具体对话或开发记录时，这一流程只能作为本组总结的协作要求，不能写成每一次编码都已经严格执行的事实。
\section{采用 AI 后的收益与不足}
AI 能帮助快速定位代码、归纳接口与提交历史、整理测试清单和排版初稿，减少重复整理工作。但它无法替成员确认真实分工、会议发生情况和未提供的使用经历，也不能把生成的解释当作已经发生的项目事实。此次文档纠正提醒我们：只有看似完整的报告还不够，必须遵循明确的保留要求，逐条核对来源。

运行时 AI 的设计边界、代码实现和问题改进分别见第三章、第四章和第九章。开发辅助 AI 的协作风险在于未经理解即接受代码、用实现逻辑反向生成测试或对生成文档缺乏核查；成员应对需求判断、资金与权限规则、事务设计、测试预期及正式提交内容负责。
